% Mean-Variance Optimization & the Efficient Frontier — research-note PDF.
%
% PROSE, CODE AND TABLES ARE VERBATIM from the website article:
%   site/src/content/articles/mvo-efficient-frontier.tsx
%   site/src/content/articles/mvo-code.ts
% Metadata from lib/articles.ts; project-card fallbacks from lib/topics.ts.
%
% CHARTS generated by figures/make_mvo_efficient_frontier.py from the article's
% own data module (data/mvo-efficient-frontier.ts) — the computed results the
% site's <Heatmap> and <Scatter> render.
\documentclass{pyportfolios-note}

\noteshorttitle{MVO \& the Efficient Frontier}

\begin{document}
\thispagestyle{notecover}

\notemasthead{Quantitative Insights}{Tutorial}{Q3 2026}{}

\notetitle{Mean-Variance Optimization\\\& the Efficient Frontier}

\notedeck{Don't pick assets — pick the combination. Building the efficient
frontier across six asset classes, and the two portfolios everyone quotes.}

\notelede{Harry Markowitz's 1952 insight launched modern portfolio theory: don't
pick assets in isolation, pick the \textit{combination} that gives the most
return per unit of risk. We build the \textbf{efficient frontier} from real data
across six asset classes and find the two portfolios everyone quotes —
\textbf{minimum variance} and \textbf{maximum Sharpe}.}

\notecard
  {Portfolio Optimization}
  {NumPy · Pandas · PyPortfolioOpt · Matplotlib}
  {SPY, TLT, GLD, VNQ, VEA, VWO}
  {Jan 2015 – Dec 2024}
  {Strategic asset allocation — the benchmark every other method is measured
   against}

% ============================================================= 1 Summary =====
\noteneedroom{150bp}
\notesection{1.\notenumsep Summary}

\begin{multicols}{2}
Mean-variance optimization treats a portfolio's expected return as the weighted
average of its assets, and its risk as the \textit{portfolio} variance — which
depends not just on each asset's volatility but on how they \textbf{co-move}.
Minimizing variance for each level of target return traces out the
\textbf{efficient frontier}: the set of portfolios you can't beat. Two points on
it get special names — the \textbf{minimum-variance} portfolio (leftmost) and
the \textbf{maximum-Sharpe} (tangency) portfolio, where the risk-adjusted return
peaks.
\end{multicols}

% =========================================================== 2 Intuition =====
\noteneedroom{170bp}
\notesection{2.\notenumsep Intuition}

\begin{multicols}{2}
\notesubsection{2.1\notenumsep The only free lunch in finance}
Combine two assets that don't move in lockstep and the portfolio's risk is
\textit{less} than the weighted average of their risks — because their wiggles
partly cancel. That cancellation is diversification, and it's why a
stock-bond-gold mix can have lower volatility than any single sleeve.

\notesubsection{2.2\notenumsep Correlation does the heavy lifting}
Volatility tells you how much one asset moves; \textbf{correlation} tells you
whether they move together. Two 15\%-vol assets with −0.3 correlation build a
far smoother portfolio than two with +0.9. The whole power of MVO lives in the
covariance matrix.

\notesubsection{2.3\notenumsep What the frontier shows}
Plot every possible portfolio in risk-return space and they fill a bullet-shaped
cloud. Only the \textbf{upper-left edge} matters — for any risk level, that's
the highest return available. Anything below it is a portfolio you'd never
rationally hold.
\end{multicols}

% ================================================= 3 Theory & Mechanics ======
\noteneedroom{180bp}
\notesection{3.\notenumsep Theory \& Mechanics}

With weights $w$, expected returns $\mu$ and covariance matrix $\Sigma$:

\[ \text{portfolio return} = w^\top \mu, \qquad
   \text{portfolio variance} = w^\top \Sigma\, w \]

The efficient frontier solves, for each target return $\mu^*$:

\[ \min_w\; w^\top \Sigma\, w \quad\text{s.t.}\quad
   w^\top \mu = \mu^*,\;\; \sum_i w_i = 1 \]

\textbf{Maximum Sharpe} maximizes $(w^\top\mu - r_f)/\sqrt{w^\top\Sigma w}$;
\textbf{minimum variance} just minimizes $w^\top\Sigma w$. We'll build the
frontier two ways — a transparent NumPy Monte-Carlo cloud, then the exact
optimum with \textbf{PyPortfolioOpt}.

% ======================================= 4 Applied Example — Six ETFs ========
\noteneedroom{200bp}
\notesection{4.\notenumsep Applied Example — Six ETFs}

\notesubsection{4.1\notenumsep A diversified multi-asset universe}

Six asset classes chosen to \textit{not} move together — the raw material of a
good frontier:

\begin{notetable}{@{}p{110bp}p{386bp}@{}}
\thcell{Ticker} & \thcell{Asset class} \\
\theadrule
SPY & US equities (S\&P 500) \\ \rowrule
TLT & Long US Treasuries \\ \rowrule
GLD & Gold \\ \rowrule
VNQ & US real estate (REITs) \\ \rowrule
VEA & Developed ex-US equities \\ \rowrule
VWO & Emerging-market equities \\
\end{notetable}

\begin{notecode}
TICKERS = ["SPY", "TLT", "GLD", "VNQ", "VEA", "VWO"]
START, END = "2015-01-01", "2024-12-31"

def load_prices(tickers, start, end):
    """Adjusted-close prices: yfinance first, Stooq as fallback."""
    try:
        import yfinance as yf
        df = yf.download(tickers, start=start, end=end, auto_adjust=True, progress=False)["Close"]
        if not df.empty:
            return df[tickers].dropna()
    except Exception as exc:
        print(f"yfinance failed ({exc}); trying Stooq…")
    cols = {}
    for t in tickers:
        url = f"https://stooq.com/q/d/l/?s={t.lower()}.us&i=d"
        s = pd.read_csv(url, parse_dates=["Date"], index_col="Date")["Close"].rename(t)
        cols[t] = s
    return pd.concat(cols, axis=1).loc[start:end].dropna()

px = load_prices(TICKERS, START, END)
rets = px.pct_change().dropna()

mu    = rets.mean() * 252          # annualised expected returns
Sigma = rets.cov() * 252           # annualised covariance
vol   = np.sqrt(np.diag(Sigma))

summary = pd.DataFrame({"ann. return": mu, "ann. vol": vol,
                        "Sharpe": mu / vol}).round(3)
print(summary)
print(f"\n{len(px)} trading days, {px.index[0].date()} → {px.index[-1].date()}")
\end{notecode}

Ten years of daily data (2,515 trading days, 2015–2024). Note how different the
standalone Sharpe ratios are — and that TLT earned essentially \textit{nothing}
over the decade:

\begin{notetable}{@{}p{136bp}>{\raggedleft\arraybackslash}p{120bp}>{\raggedleft\arraybackslash}p{120bp}>{\raggedleft\arraybackslash}p{100bp}@{}}
\thcell{Ticker} & \thcell{Ann. return} & \thcell{Ann. vol} & \thcell{Sharpe} \\
\theadrule
SPY & 13.9\% & 17.6\% & 0.79 \\ \rowrule
TLT & −0.0\% & 15.3\% & −0.00 \\ \rowrule
GLD & 8.5\% & 14.1\% & 0.60 \\ \rowrule
VNQ & 6.9\% & 20.8\% & 0.33 \\ \rowrule
VEA & 6.9\% & 17.3\% & 0.40 \\ \rowrule
VWO & 6.0\% & 19.8\% & 0.30 \\
\end{notetable}

\notesubsection{4.2\notenumsep The correlation matrix — the source of the free lunch}

Note where correlations are \textit{low}: Treasuries (TLT, −0.21 to SPY) and
gold (GLD, +0.05 to SPY) barely track equities, which is exactly why they smooth
the portfolio.

\begin{notecode}
corr = rets.corr()

fig, ax = plt.subplots(figsize=(6.5, 5.5))
im = ax.imshow(corr, cmap="coolwarm", vmin=-1, vmax=1)
ax.set_xticks(range(len(TICKERS))); ax.set_xticklabels(TICKERS)
ax.set_yticks(range(len(TICKERS))); ax.set_yticklabels(TICKERS)
for i in range(len(TICKERS)):
    for j in range(len(TICKERS)):
        ax.text(j, i, f"{corr.iloc[i,j]:.2f}", ha="center", va="center",
                color="white" if abs(corr.iloc[i,j]) > 0.5 else "black", fontsize=9)
fig.colorbar(im, fraction=0.046, pad=0.04)
ax.set_title("Correlation of daily returns")
plt.tight_layout(); plt.show()
\end{notecode}

\notefigure{figures/mvo-efficient-frontier/fig1-corr.png}
  {Figure 4.2 · Correlation of daily returns, 2015–2024}
  {The equity block (SPY, VNQ, VEA, VWO) runs 0.55 to 0.86; TLT is slightly
   negative against equities and GLD is near zero}
  {Source: Author calculations, from the article's computed results.}

\notesubsection{4.3\notenumsep The efficient frontier — Monte-Carlo cloud}

Generate 20,000 random portfolios to \textit{see} the bullet, then the frontier
as its upper-left edge. Colour = Sharpe ratio.

\begin{notecode}
N = 20_000
n = len(TICKERS)
rf = 0.02

w = np.random.dirichlet(np.ones(n), N)          # random long-only weights summing to 1
port_ret = w @ mu.values
port_vol = np.sqrt(np.einsum("ij,jk,ik->i", w, Sigma.values, w))
port_sharpe = (port_ret - rf) / port_vol

fig, ax = plt.subplots(figsize=(10, 6))
sc = ax.scatter(port_vol, port_ret, c=port_sharpe, cmap="viridis", s=6, alpha=0.5)
# individual assets
ax.scatter(vol, mu.values, c="red", marker="D", s=60, zorder=5)
for i, t in enumerate(TICKERS):
    ax.annotate(t, (vol[i], mu.values[i]), textcoords="offset points", xytext=(7, 3))
fig.colorbar(sc, label="Sharpe ratio")
ax.set_xlabel("Annualised volatility"); ax.set_ylabel("Annualised return")
ax.set_title(f"{N:,} random portfolios — the efficient frontier bullet")
plt.tight_layout(); plt.show()
\end{notecode}

\notesubsection{4.4\notenumsep The exact optima with PyPortfolioOpt}

The cloud shows the shape; PyPortfolioOpt solves for the \textit{exact}
minimum-variance and maximum-Sharpe portfolios and overlays the true frontier.

\begin{notecode}
from pypfopt import EfficientFrontier, expected_returns, risk_models, plotting

mu_pp = expected_returns.mean_historical_return(px)
S_pp  = risk_models.sample_cov(px)

# max Sharpe
ef = EfficientFrontier(mu_pp, S_pp)
ef.max_sharpe(risk_free_rate=rf)
w_sharpe = ef.clean_weights()
perf_sharpe = ef.portfolio_performance(risk_free_rate=rf)

# min variance
ef2 = EfficientFrontier(mu_pp, S_pp)
ef2.min_volatility()
w_minvar = ef2.clean_weights()
perf_minvar = ef2.portfolio_performance(risk_free_rate=rf)

weights = pd.DataFrame({"Max Sharpe": w_sharpe, "Min Variance": w_minvar}).round(3)
print(weights)
print(f"\nMax Sharpe : return {perf_sharpe[0]:.1%}, vol {perf_sharpe[1]:.1%}, Sharpe {perf_sharpe[2]:.2f}")
print(f"Min Variance: return {perf_minvar[0]:.1%}, vol {perf_minvar[1]:.1%}, Sharpe {perf_minvar[2]:.2f}")
\end{notecode}

\begin{notetable}{@{}p{180bp}>{\raggedleft\arraybackslash}p{150bp}>{\raggedleft\arraybackslash}p{150bp}@{}}
\thcell{Ticker} & \thcell{Max Sharpe} & \thcell{Min Variance} \\
\theadrule
SPY & 0.563 & 0.268 \\ \rowrule
TLT & 0.000 & 0.371 \\ \rowrule
GLD & 0.437 & 0.280 \\ \rowrule
VNQ & 0.000 & 0.000 \\ \rowrule
VEA & 0.000 & 0.081 \\ \rowrule
VWO & 0.000 & 0.000 \\
\end{notetable}

\textbf{Max Sharpe}: 10.8\% return at 11.9\% vol (Sharpe 0.73) — a two-asset
SPY + GLD barbell. \textbf{Min variance}: 5.7\% return at 9.3\% vol (Sharpe
0.40) — note it holds 37\% TLT \textit{despite} TLT's zero return, purely for its
negative correlation. That is the free lunch in action: the optimizer pays for
co-movement, not for standalone performance.

\notefigure{figures/mvo-efficient-frontier/fig2-frontier.png}
  {Figure 4.4 · 20,000 random portfolios, the exact frontier, and the two special portfolios}
  {The random portfolios form a bullet; the frontier runs along its upper-left
   edge, with the capital market line tangent at maximum Sharpe}
  {Source: Author calculations, from the article's computed results. The cloud is
   drawn from the stored sample of the simulation.}

Every single ETF plots \textit{inside} the cloud, well below the frontier — even
SPY, the decade's best performer, sits under the line. No individual asset is
efficient; only combinations are.

% ========================================================== 5 Conclusion =====
\noteneedroom{330bp}
\notesection{5.\notenumsep Conclusion}

\begin{multicols}{2}
\notesubsection{5a.\notenumsep Strengths}
\begin{notebullets}
\item \textbf{Quantifies diversification} — turns ``don't put all your eggs in
      one basket'' into an exact weight vector
\item \textbf{One framework, any universe} — stocks, bonds, gold, real estate
      all go in the same optimizer
\item \textbf{Closed-form and fast} — the frontier is a quadratic program that
      solves instantly
\item \textbf{The foundation} — every allocation method (risk parity,
      Black-Litterman, CVaR) is a response to MVO
\end{notebullets}

\notesubsection{5b.\notenumsep Weaknesses \& Limitations}
\begin{notebullets}
\item \textbf{Garbage in, garbage out} — expected returns are notoriously hard
      to estimate, and MVO is hypersensitive to them
\item \textbf{Concentrated, unstable weights} — tiny input changes can swing
      allocations wildly (the ``error-maximization'' critique)
\item \textbf{Backward-looking} — historical covariance assumes the past regime
      persists
\item \textbf{Ignores tail risk} — variance treats upside and downside
      symmetrically; crashes aren't Gaussian
\end{notebullets}

\notesubsection{5c.\notenumsep Applications in Practice}
\begin{notebullets}
\item Strategic asset allocation for pension funds and endowments
\item Setting the neutral portfolio a discretionary manager tilts away from
\item The benchmark every alternative allocation method is measured against
\end{notebullets}

\notesubsection{5d.\notenumsep Alternatives \& Extensions}
\begin{notebullets}
\item \textbf{Black-Litterman} — fixes the input-sensitivity problem by blending
      market equilibrium with views (the next tutorial)
\item \textbf{Risk parity} — sidesteps return estimation entirely by allocating
      on risk contribution (our futures-based piece)
\item \textbf{Hierarchical Risk Parity} — uses clustering instead of matrix
      inversion for stabler weights
\item \textbf{CVaR optimization} — replaces variance with a genuine tail-risk
      measure
\end{notebullets}
\end{multicols}

\end{document}
